%% ==========================================================================
%% Hauptteil der Bachelorarbeit
%% ==========================================================================

\section{Hauptteil}
\label{sec:hauptteil}

\lipsum[1]

\subsection{Theoretische Grundlagen}
\label{subsec:theorie}

Wissenschaftliche Arbeiten erfordern eine fundierte Auseinandersetzung mit dem aktuellen Forschungsstand. Nach \textcite{blinowski_monolithic_2022} lassen sich verschiedene Architekturansätze bezüglich ihrer Performanz und Skalierbarkeit differenziert evaluieren. Auch grundlegende methodische Ansätze können über direkte und indirekte Zitate eingebunden werden \parencite[S.~42]{nussbaum}.

\subsubsection{Konzepte und Begriffsabgrenzung}
\label{subsubsec:konzepte}

Gemäß den IU-Richtlinien können Arbeiten in maximal drei Gliederungsstufen unterteilt werden. Hierbei erfolgt die Formatierung der Überschriften der 3.~Stufe in 11\,Pt. fett. 

\lipsum[2]

\subsection{Methodik und empirische Untersuchung}
\label{subsec:methodik}

Die empirische Erhebung erfolgte anhand leitfadengestützter Experteninterviews. Der vollständige Erhebungsbogen ist im Anhang dokumentiert (siehe Anhang~A).

In Abbildung~\ref{fig:logo_iu} ist beispielhaft das Logo der \ac{iu} dargestellt. Abbildungen und Tabellen sollten stets im Fließtext referenziert und mit einer eindeutigen Quellenangabe in Schriftgröße 10\,Pt. versehen werden.

\begin{figure}[htbp]
    \centering
    \includegraphics[width=0.45\textwidth]{logos/iu_Logo_EN.jpg}
    \caption[Logo der \acs{iu}]{Logo der \ac{iu}}
    \label{fig:logo_iu}
    \quelle{In Anlehnung an \url{https://www.iu.de}}
\end{figure}

Zur Veranschaulichung von strukturierten Daten dient Tabelle~\ref{tab:beispiel}. Laut IU-Richtlinien ist ein Tabellenverzeichnis ab drei verwendeten Tabellen verpflichtend.

\begin{table}[htbp]
    \centering
    \caption{Übersicht exemplarischer Bewertungskriterien}
    \label{tab:beispiel}
    \vspace{4pt}
    \begin{tabular}{llr}
        \hline
        \textbf{Kriterium} & \textbf{Beschreibung} & \textbf{Gewichtung} \\
        \hline
        Inhaltliche Tiefe & Systematische Bearbeitung der Forschungsfrage & 40\,\% \\
        Wissenschaftlichkeit & Literaturarbeit nach APA-Richtlinien & 30\,\% \\
        Struktur \& Formalia & Einhaltung des IU-Leitfadens und Layouts & 30\,\% \\
        \hline
    \end{tabular}
    \quelle{Eigene Darstellung}
\end{table}

\subsection{Quellcode-Darstellung}
\label{subsec:code}

Quellcode kann übersichtlich über das \texttt{listings}-Paket eingebunden werden:

\begin{lstlisting}[language=Java, caption={Beispielhafte Java-Methode}, label={lst:java_beispiel}]
public class Application {
    public static void main(String[] args) {
        System.out.println("Willkommen an der IU!");
    }
}
\end{lstlisting}

\lipsum[7]
\lipsum[7-8]